\clearpage\begingroup\HistoricalHead{223}{R}{BESTIMMUNG DER BAHNEN DER HIMMELSKÖRPER.}
\makeatletter\tagsleft@true\let\veqno\@@leqno\makeatother
\noindent{\small [S. 3]}
\begin{center}
[3.]\par\medskip
\textls[100]{Vorschriften zur Berechnung der Elemente, aus zweien geocentrischen Örtern, der Zwischenzeit, und den zugehörigen Abständen}\textup{[\textsuperscript{**})]}.
\end{center}
\begingroup\renewcommand{\thefootnote}{}\footnotetext{[\textsuperscript{**}) Vgl. Theoria motus, art. 124, Werke VII, 1906, S. 165.]}\endgroup
\[
\rho=\surd\{RR+\Delta\Delta-2R\Delta\cos(V-\alpha)\}.
\]
\noindent Ganz allgemein
\[
\left\{\begin{aligned}
\Delta\sin\alpha-R\sin V&=\rho\sin\lambda\\
\Delta\cos\alpha-R\cos V&=\rho\cos\lambda
\end{aligned}\right\};
\]
\noindent folglich
\begin{equation}
\left\{\begin{aligned}
\Delta\sin(V-\alpha)&=\rho\sin(V-\lambda)\\
\Delta\cos(V-\alpha)-R&=\rho\cos(V-\lambda)
\end{aligned}\right\},\tag*{I.}
\end{equation}
\begin{equation}
\left\{\begin{aligned}
R\sin(V-\alpha)&=\rho\sin(\alpha-\lambda)\\
R\cos(V-\alpha)-\Delta&=-\rho\cos(\alpha-\lambda)
\end{aligned}\right\},\tag*{II.}
\end{equation}
\begin{equation}
\left\{\begin{aligned}
(\Delta+R)\sin\tfrac12(V-\alpha)&=\rho\sin\{\tfrac12(V+\alpha)-\lambda\}\\
(\Delta-R)\cos\tfrac12(V-\alpha)&=\rho\cos\{\tfrac12(V+\alpha)-\lambda\}
\end{aligned}\right\}\tag*{III.}
\end{equation}
\clearpage
\HistoricalHead{224}{L}{THEORETISCHE ASTRONOMIE. NACHLASS.}\SourceLeaf{229}
\[
r=\rho\sec\beta,\qquad
\vartheta=\operatorname{tang}\beta=\frac{\Delta}{\rho}\operatorname{tang}\delta=\frac{\Delta\Theta}{\rho}.
\]
\noindent Allgemein [$\vartheta=g\sin(\lambda-\AstNode)$, also]
\[
\left\{\begin{aligned}
\vartheta'\sin\lambda-\vartheta\sin\lambda'&=g\sin(\lambda'-\lambda)\sin\AstNode\\
\vartheta'\cos\lambda-\vartheta\cos\lambda'&=g\sin(\lambda'-\lambda)\cos\AstNode
\end{aligned}\right\}
\]
\[
g\sin(\lambda'-\lambda)=\surd\bigl(\vartheta\vartheta+\vartheta'\vartheta'-2\vartheta\vartheta'\cos(\lambda'-\lambda),
\]
\noindent folglich
\begin{equation}
\left\{\begin{aligned}
\vartheta\sin(\lambda'-\lambda)&=g\sin(\lambda'-\lambda)\sin(\lambda-\AstNode)\\
\vartheta'-\vartheta\cos(\lambda'-\lambda)&=g\sin(\lambda'-\lambda)\cos(\lambda-\AstNode)
\end{aligned}\right\},\tag*{I.}
\end{equation}
\begin{equation}
\left\{\begin{aligned}
\vartheta'\sin(\lambda'-\lambda)&=g\sin(\lambda'-\lambda)\sin(\lambda'-\AstNode)\\
\vartheta'\cos(\lambda'-\lambda)-\vartheta&=g\sin(\lambda'-\lambda)\cos(\lambda'-\AstNode)
\end{aligned}\right\},\tag*{II.}
\end{equation}
\begin{equation}
\begin{aligned}
(\vartheta'+\vartheta):2\cos\tfrac12(\lambda'-\lambda)&=g\sin\bigl(\tfrac12(\lambda+\lambda')-\AstNode\bigr),\\
(\vartheta'-\vartheta):2\sin\tfrac12(\lambda'-\lambda)&=g\cos\bigl(\tfrac12(\lambda+\lambda')-\AstNode\bigr),
\end{aligned}\tag*{III.}
\end{equation}
\[
\begin{aligned}
\operatorname{tang}(v-\AstNode)&=\frac{\operatorname{tang}(\lambda-\AstNode)}{\cos i},\\
\cos(v-\AstNode)&=\cos(\lambda-\AstNode)\cos\beta,\\
\sin(v-\AstNode)&=\frac{\sin\beta}{\sin i},
\end{aligned}
\]
\[
\operatorname{tang}(v'-v)=\sec i\mathbin{.}\frac{\sin(\lambda'-\lambda)}{\cos(\lambda'-\lambda)+\vartheta\vartheta'}.
\]

Hiedurch hat man also die beiden Örter in der Bahn und zugleich die wahren Abstände von der Sonne [sowie die Länge des aufsteigenden Knotens und die Neigung der Bahn].
\clearpage\endgroup
